\paragraph{First-order equivariant class on the Heisenberg background with curvature: Landau sectors of the linearised equations (conditional; not closure).}
Take the wider ansatz \(w_a{}^{cd}=k_a(\chi B^{cd}+\zeta N^{cd}+xN'^{cd})\) of the preceding paragraph on time times the Heisenberg group, with its normalisation and the fluctuation \(\mathbb D_a\to\mathbb D_a-h_a{}^c\mathbb D_c+\tfrac12\omega_a{}^{cd}\mathcal S_{cd}\), in the sectors of the paragraph on left-invariant backgrounds with torsion: \(E_0=\mathrm iq_0\), \(E_2=\mathrm ip\) with \(p>0\), and \(E_1\), \(E_3\) acting on oscillator levels. The linearised equations have real coefficients, so the sectors with \(p<0\) follow by complex conjugation. For a constant connection \(\mathsf F_{ab}[A]=E_aA_b-E_bA_a-\mathsf c_{ab}{}^dA_d\) contains no \(w\), and on this family
\[
\mathcal Q_b[A]=\mathcal Q_b[A]\big|_{w=0}-\chi\,k^a\mathsf F_{ab}[A]
\]
for all \(\zeta,x\) and all \(A\) (exact), so the connection adds only a gauge-invariant term of first order along the null \(k\). On \(A=(0,\phi,0,\mathrm i\phi)\), with \(\phi\) as in that paragraph, this gives \(\mathcal Q_b[A]=\bigl(q_0^2-p^2+\kappa p+\mathrm i\chi(q_0-p)\bigr)A_b\) (exact). The operator \(\mathcal Q\) agrees with the linearised third bracket of the operators on the bundle, order-zero part included (exact at a sampled point for polynomial \(A\), at two parameter sets). With \(q_\parallel=(q_0,0,p,0)\), every characteristic root of the forty equations in these sectors, a value of \(q_0\) at which their rank drops below its generic value, satisfies
\[
\bigl(q_\parallel+\tfrac{\mathrm i}2c\chi k\bigr)^2+\kappa p\,l=0,\quad\text{that is}\quad q_0^2-p^2+\mathrm ic\chi(q_0-p)=\kappa p\,l,\qquad c\in\{-1,1,3\},\ l\in\{-1,1,3,\dots\}.
\]
For \(l\ge3\) the numbers of these quadratic factors are \(10\), \(6\) and \(4\) for \(c=1,3,-1\), twenty in all as at \(\chi=0\), and half of these at \(l=\pm1\); \(\mathcal Q\) has \(c=1\), and the class \(c=1\) has five times its multiplicities (numerical, at sampled \(\kappa,\chi,p\)). At \(\zeta=x=0\) a rotation of the \(1\)-\(3\) plane combined with the oscillator level is conserved, and in each of its blocks \(J=-2,\dots,3\), which contain all factors with \(l\le5\), the greatest common divisor of two random compressions of the maximal minors, which every characteristic root annuls, is exactly the product of these factors at \((\kappa,p,\chi)=(1,\tfrac12,\tfrac12)\) and \((2,\tfrac14,-\tfrac13)\) (exact in rational arithmetic). For \(\zeta,x\neq0\) the same values and counts appear for \(l\le5\), and in the sector \(p=0\), where \(E_1\) and \(E_3\) commute, the roots lie on the complex cones \((q+\tfrac{\mathrm i}2c\chi k)^2=0\) with the same \(c\) (numerical, at sampled points). The kernel at generic \(q_0\) is the ten gauge directions for each level of the parameter, and the frame labels along \(E_0\) and \(E_2\) reduce to \(\mathcal Q\) as on family~(i) (numerical). At \(\chi=0\) this gives the levels \(q_0^2=p^2+\kappa p\,l\) of the paragraph on left-invariant backgrounds with torsion for every \(\zeta,x\) (numerical).

So \(\chi\), \(\zeta\) and \(x\) shift no level: the values \(\kappa p\,l\) are those at \(\chi=0\), and \(\chi\) enters only through the imaginary null vector \(\tfrac{\mathrm i}2c\chi k\) added to \(q_\parallel\), whose contribution \(\mathrm ic\chi\,k^aq_a\) is linear in the momentum. No gauge-invariant gap or mass therefore appears for the Lorentz or torsion modes at this order, although the background carries the scale \(\kappa\) and \(\mathsf R_{abcd}\mathsf R^{abcd}=-4\chi^2\kappa^2\neq0\), and no sign of a mass arises. The roots are \(q_0=-\tfrac{\mathrm i}2c\chi\pm\bigl((p+\tfrac{\mathrm i}2c\chi)^2+\kappa p\,l\bigr)^{1/2}\). For \(\chi\neq0\) none is real at \(p>0\), and their imaginary parts have both signs; for \(l\ge1\), \(|\operatorname{Im}q_0|<|c\chi|\), with supremum \(|c\chi|\) as \(p\to\infty\); for \(l=-1\) the largest value of \(|\operatorname{Im}q_0|\) is \(\tfrac12|c\chi|+\sigma_c\) with \(\sigma_c=\tfrac\kappa4+(\tfrac{\kappa^2}{16}+\tfrac14c^2\chi^2)^{1/2}\), attained at \(p=\sigma_c\) (by hand, checked numerically). For \(c=1\) this is the mode \(A\) above: the growing mode of the paragraph on left-invariant backgrounds with torsion persists at \(\chi\neq0\), its largest rate rises from \(\kappa/2\) to \(\tfrac12|\chi|+\sigma_1>\kappa/2\), and the growth is no longer confined to \(0<p<\kappa\). Over all modes the largest rate is \(\tfrac32|\chi|+\sigma_3\). These are section values and no norm is used, so, as for family~(i), no statement on stability is made.

These statements are for Equation~\eqref{eq:Smatrix} with \(C_{\mu\nu\lambda}=0\) and no matter term, tree level, linear order in the fluctuation, the Lorentzian contraction and the wider family on time times the Heisenberg group with constant connection, in the sectors named. A norm, the nonlinear evolution of the growing modes, loops, for which the cancellation at second order about commuting matrices does not apply since the commutators of this background are nonzero, the solutions with \(\mathsf R_{abcd}\mathsf R^{abcd}\neq0\) of the survey of the preceding paragraph, the regulated trace and the domains of the unbounded operators are not treated. On the Heisenberg algebra, constant connections with \(\eta^{ab}w_a{}^{cd}w_b{}^{ef}\neq0\), a form that vanishes on the wider family since \(k\) is null, also solve all three brackets, most of them with \(\mathsf R_{abcd}\mathsf R^{abcd}>0\) (numerical survey, not a classification); they are not linearised here. Nothing here shows that the matrix integral selects this class, these backgrounds or the algebra \(\mathcal A\), or yields the frame and its non-degeneracy, so Assumptions~\ref{ass:frame-lock}, \ref{ass:nondeg}, \ref{ass:lorentzian}, \ref{ass:local-lorentz} and~\ref{ass:EFT-a2} are not discharged, and the frame of these backgrounds is not identified with the emergent metric of Chapter~\ref{ch:vierbein-emergence}; Problems~D, F and~N remain open and PROGRAMME-CONDITIONAL.
